\chapter{Fórmulas, unidades y comprobaciones}\label{ap:formulas}

Una fórmula solo es útil si se conocen sus unidades y el problema para el que fue derivada.
Este apéndice reúne las relaciones principales y añade controles que ayudan a detectar
errores antes de que lleguen a una gráfica. Los capítulos contienen su desarrollo.

\section{Constantes y conversiones}
\begin{center}
\renewcommand{\arraystretch}{1.2}
\begin{tabular}{ll}
\toprule
Magnitud & Valor utilizado \\
\midrule
Velocidad de la luz $c$ & \SI{299792458}{\metre\per\second} \\
Constante de Planck $h$ & \SI{6.62607015e-34}{\joule\second} \\
Constante de Boltzmann $k_{\mathrm B}$ & \SI{1.380649e-23}{\joule\per\kelvin} \\
Stefan--Boltzmann $\sigma$ & $\SI{5.670374419e-8}{\watt\per\square\metre\per\kelvin\tothe{4}}$ \\
Un micrómetro & $10^{-6}$ m \\
Un megavatio & $10^6$ W $=10^6$ J/s \\
Una hectárea & $10^4$ m$^2$ \\
Un kilómetro cuadrado & $10^6$ m$^2=100$ ha \\
Temperatura absoluta & $T\,[\mathrm K]=t\,[{}^\circ\mathrm C]+273{,}15$ \\
\bottomrule
\end{tabular}
\end{center}
Las tres primeras constantes tienen valores fijados en el SI; las cifras de $\sigma$
se muestran redondeadas. En los ejemplos basta normalmente usar $5{,}67\times10^{-8}$.
Nunca se sustituye una temperatura en Celsius en $T^4$ ni en la ley de Planck.

\section{Radiación térmica}

Para un cuerpo negro, la radiancia espectral por unidad de longitud de onda es
\begin{equation}
 B_\lambda(T)=\frac{2hc^2}{\lambda^5}
 \frac{1}{\exp(hc/(\lambda k_{\mathrm B}T))-1}.
\end{equation}
Con $\lambda$ en metros, el resultado está en
$\mathrm{W\,m^{-2}\,sr^{-1}\,m^{-1}}$. Para expresarlo por micrómetro se multiplica por
$10^{-6}$. Esa conversión no es lo mismo que reemplazar sin más la longitud de onda
por su valor numérico en micrómetros.

La forma práctica del capítulo~\ref{cap:radiacion} usa
\begin{equation}
 B_\lambda(T)=\frac{c_1}{\lambda^5[\exp(c_2/(\lambda T))-1]},\quad
 c_1\simeq1{,}19104\times10^8,\quad c_2\simeq14387{,}77,
\end{equation}
con $\lambda$ en micrómetros y radiancia en
$\mathrm{W\,m^{-2}\,sr^{-1}\,\mu m^{-1}}$. Su inversión da
\begin{equation}
 T_b=\frac{c_2}{\lambda\ln[1+c_1/(\lambda^5L_\lambda)]}.
\end{equation}
El control más simple consiste en generar $L=B_\lambda(T)$ y comprobar que la inversión
recupera $T$. En un instrumento real se integra sobre la respuesta espectral del canal;
la longitud central usada aquí es una aproximación pedagógica.

\begin{align}
 \lambda_{\max}T&\simeq2898\ \mathrm{\mu m\,K},\\
 M&=\sigma T^4,\qquad P=A\varepsilon\sigma T^4.
\end{align}
Wien localiza el máximo de $B_\lambda$, no el de una curva expresada por frecuencia.
$M$ es potencia por área integrada en el hemisferio. Para emisión isotrópica, integrar
la radiancia sobre el hemisferio incluye el factor geométrico $\pi$; no se deben
intercambiar $B_\lambda$ y $M$ como si fueran la misma magnitud.\fuente{openstax}

\section{Píxel mixto y potencia del fuego}

Con emisividad uno, sin atmósfera ni reflexión, y fuego que ocupa una fracción $p$,
\begin{equation}
 L_\lambda=pB_\lambda(T_f)+(1-p)B_\lambda(T_b).
\end{equation}
Las \emph{radiancias} se promedian; las temperaturas de brillo no. La inversión de
Planck es no lineal. Para una superficie de fuego $A_f$ a temperatura uniforme,
\begin{equation}
 P_{\mathrm{exceso}}\simeq A_f\varepsilon\sigma(T_f^4-T_b^4).
\end{equation}
El modelo es idealizado: un fuego real tiene una distribución de temperaturas.

La aproximación MIR del capítulo~\ref{cap:frp} relaciona un exceso de radiancia con
potencia mediante un coeficiente que depende de banda y rango de temperaturas:
\begin{equation}
 P\simeq\frac{\sigma}{a}A_{\mathrm{pixel}}
 [L_{\mathrm{MIR}}-L_{\mathrm{fondo}}].
\end{equation}
Se debe usar $a$ con las mismas unidades espectrales de $L$. No es una constante
universal de todos los sensores. Para usar FIRMS normalmente se parte de la FRP
publicada, siguiendo la documentación del producto.\fuente{wooster2003}

Si se conoce una curva temporal de potencia,
\begin{equation}
 E=\int_{t_0}^{t_1}P(t)\,\mathrm dt
 \simeq\sum_i\frac{P_i+P_{i+1}}{2}(t_{i+1}-t_i).
\end{equation}
Con MW y segundos se obtiene MJ. Una potencia de 10 MW mantenida durante una hora
equivale a $10\times3600=36000$ MJ, es decir, 36 GJ. Dos pasadas con FRP de 10 MW no
demuestran que la potencia permaneciera constante entre ellas. La integral exige
un modelo temporal o muestreo suficiente.\fuente{wooster2005}

\section{Órbita, geometría y coordenadas}

En una órbita circular de radio $r=R+h$, con parámetro gravitacional $\mu=GM$,
\begin{equation}
 T_{\mathrm{orb}}=2\pi\sqrt{\frac{r^3}{\mu}}.
\end{equation}
Si $r$ está en metros y $\mu$ en $\mathrm{m^3\,s^{-2}}$, el periodo sale en segundos.
El periodo de una vuelta no es el intervalo entre observaciones de un punto del suelo.

Para una esfera de radio $R$, una celda entre latitudes $\phi_1$ y $\phi_2$ y ancho
longitudinal $\Delta\lambda$ tiene área
\begin{equation}\label{eq:area-celda}
 A=R^2\Delta\lambda(\sin\phi_2-\sin\phi_1),
\end{equation}
con ángulos en radianes. La aproximación local es
$A\simeq R^2\Delta\phi\Delta\lambda\cos\phi$. Esto explica por qué una rejilla
regular en grados no tiene área constante. En un elipsoide se usa una formulación
geodésica o una proyección adecuada si esa precisión es necesaria.

Los índices de una rejilla de paso $\Delta$ en grados son
\begin{equation}
 i=\left\lfloor\frac{\phi+90}{\Delta}\right\rfloor,\qquad
 j=\left\lfloor\frac{\lambda+180}{\Delta}\right\rfloor.
\end{equation}
La operación es suelo matemático, no truncamiento hacia cero. Se define además la
convención de bordes, por ejemplo intervalos cerrados a la izquierda y abiertos a la
derecha. Los polos y $180^\circ$ necesitan tratamiento de frontera. En el caso peruano
se usa aritmética decimal para no desplazar un punto exactamente sobre un borde por
un error binario de redondeo.

\section{Conteos y calibración}

Bajo el modelo Poisson, $\operatorname{Var}(N)=\mu$; estimar $\mu$ por $N$ lleva a la
aproximación relativa $1/\sqrt N$ para $N>0$. No es una ley universal de incertidumbre
del fuego. Dependencias, omisiones y muestreo desigual pueden dominarla. Con $N=0$
no se concluye incertidumbre cero ni se evalúa esa expresión.

Para parejas $(X_i,Y_i)$ agregadas sobre soportes comparables, la razón de sumas es
\begin{equation}
 \widehat r=\frac{\sum_iY_i}{\sum_iX_i},\qquad\sum_iX_i>0.
\end{equation}
Promediar razones mensuales sin pesos estima otra cantidad. Además, que un mes con
$X_i=0$ no impida calcular la razón global no significa que el modelo pueda recuperar
su actividad: $\widehat r\,0=0$ aunque VIIRS haya registrado detecciones.

Para la contracción del capítulo~\ref{cap:armonizacion},
\begin{equation}
 r_s=\frac{Y_s+kr_g}{X_s+k}
 =\frac{X_s}{X_s+k}\frac{Y_s}{X_s}
 +\frac{k}{X_s+k}r_g\quad(X_s>0).
\end{equation}
La segunda forma muestra los pesos. $k$ tiene unidades de conteo de referencia y no
equivale necesariamente a $k$ observaciones independientes. Debe elegirse usando solo
datos de ajuste y validación interna.

Para el producto de dos factores, la propagación de primer orden da
\begin{equation}
 u_{12}^2\simeq u_1^2+u_2^2+2\rho u_1u_2,
\end{equation}
donde $u_i$ son incertidumbres estándar relativas y $\rho$ su correlación. La suma en
cuadratura supone $\rho=0$. No se introduce directamente el ancho de un intervalo
del 95\,{}\% como si fuera una desviación estándar. Cuando los factores comparten
años o sensores, un remuestreo conjunto conserva mejor sus dependencias.

\section{Métricas del calendario y de validación}

Con referencia $b_m>0$, la anomalía relativa es $100(C_m/b_m-1)$; con $b_m=0$ se
informa la diferencia absoluta y que la razón es indefinida. El perfil estacional
normalizado por año cumple
\begin{equation}
 \sum_m p_{y,m}=1,\qquad \sum_m\bar p_m=1.
\end{equation}
Estas identidades detectan meses duplicados o denominadores mal construidos.

Para $n$ predicciones de prueba y residuos $e_i=\widehat y_i-y_i$,
\begin{align}
 \mathrm{sesgo}&=\frac1n\sum_i e_i, &
 \mathrm{MAE}&=\frac1n\sum_i|e_i|,\\
 \mathrm{RMSE}&=\sqrt{\frac1n\sum_i e_i^2}, &
 \mathrm{WAPE}&=100\frac{\sum_i|e_i|}{\sum_i y_i}.
\end{align}
WAPE requiere un total observado positivo; MAPE requiere además observaciones
individuales no nulas. RMSE da mayor peso a errores grandes. Ninguna métrica reemplaza
la inspección del error por mes, zona y calidad de observación.

\begin{ejemplo}{un cero en la referencia observada}
Para valores observados $0,10,50,100$ y predicciones $2,8,45,110$, calcula MAE y WAPE.
\solucion Los errores absolutos son $2,2,5,10$; suman 19. Por tanto,
$\mathrm{MAE}=19/4=4{,}75$ y $\mathrm{WAPE}=100(19/160)=11{,}875\,{}\%$.
MAPE no está definido porque el primer observado es cero. Eliminar ese caso sin
explicarlo escondería una falsa detección de actividad.
\end{ejemplo}

\section{Resultados que sirven para revisar el libro}

El programa \texttt{scripts/verificar\_ejemplos.py} reproduce ejemplos aritméticos
didácticos. El programa \texttt{scripts/analizar\_peru.py} calcula por separado el caso
real; sus cifras no se mezclan con las inventadas para explicar una operación.

\begin{center}
\begin{tabular}{lr}
\toprule
Comprobación del caso real & Resultado \\
\midrule
Filas en los 24 CSV originales de Perú & \PeruFilas \\
Celdas-día Aqua, ajuste 2018--2020 & \PeruAjusteAqua \\
Celdas-día VIIRS, ajuste 2018--2020 & \PeruAjusteVIIRS \\
Razón de sumas del ajuste & \PeruRatio \\
WAPE mensual, prueba 2021--2022 & \PeruWAPE\,{}\% \\
Mediana de septiembre, VIIRS 2013--2023 & \PeruSeptMediana \\
Septiembre de 2024, misma selección & \PeruSeptActual \\
\bottomrule
\end{tabular}
\end{center}
Si otra ejecución difiere, primero se comparan huellas de archivos, versiones,
filtros, zona horaria y claves de agrupación. Solo después se investiga la aritmética.
